\section{Demostraciones complementarias de SP0}
\label{app:sp0-proofs}

SP0 relaciona la complejidad de la asignación uno a uno con la dinámica de mejora del juego. El resultado de complejidad usa costes normalizados. La prueba de terminación opera sobre un espacio de acciones finito, y la factibilidad del equilibrio invoca \(\lambda>\kappa_{\max}\).

\subsection{Integralidad y complejidad}
\label{app:proof-sp0-complexity}

La matriz de restricciones de~\eqref{eq:sp0-oracle} es, salvo signos, la matriz de incidencia nodo--arista del grafo bipartito robot--tarea y es totalmente unimodular. Como el lado derecho es entero, los vértices de la relajación lineal son enteros. Añadir $N-K$ tareas ficticias convierte cada solución rectangular en una permutación cuadrada sin cambiar el coste real. Para $N\geq K$, $P(N,K)=N!/(N-K)!$ denota el número de permutaciones parciales de $K$ robots tomados de una flota de $N$. El algoritmo húngaro resuelve el problema denso en $O(N^3)$ \parencite{kuhn1955Hungarian}. Una enumeración directa cuesta $\Theta(KP(N,K))$.

La selección opcional de cargas introduce una frontera distinta. Sea $y_k\in\{0, 1\}$ la decisión de servir la carga $k$, con $\sum_i x_{ik}=n_ky_k$ y $\sum_kx_{ik}\leq1$. Si los robots son intercambiables y los costes nulos, $y$ es factible si y solo si $\sum_kn_ky_k\leq N$. Maximizar $\sum_kv_ky_k$ coincide con mochila 0--1: $n_k$ es el peso, $v_k$ el valor y $N$ la capacidad. La reducción es polinómica y establece NP-dificultad \parencite{karp1972reducibility}.

\subsection{Potencial, Nash y terminación}

\paragraph{Potencial exacto.}
Para una desviación unilateral $a_i\rightarrow b_i$, el perfil que retira al robot $i$ es idéntico antes y después. Por la utilidad marginal,
\[
u_i(b_i,a_{-i})-u_i(a_i,a_{-i})
=\Phi_\lambda(b_i,a_{-i})-\Phi_\lambda(a_i,a_{-i}),
\]
que satisface la definición de potencial exacto \parencite{monderer1996potential}.

\paragraph{Caracterización del equilibrio.}
Sea $a$ infactible. Si falta una tarea y existe un robot inactivo, asignarlo reduce $D$ en uno y aumenta $C$ a lo sumo en $\kappa_{\max}<\lambda$. Si no hay robots inactivos, $N\geq K$ y una tarea está vacía, otra está duplicada. Trasladar un robot reduce simultáneamente $D$ y $E$, con variación de coste no mayor que $\kappa_{\max}$. Si solo existe exceso, retirar un robot redundante reduce $E$ y no aumenta $C$. Todo perfil infactible admite así una mejora estricta.

En un perfil factible, retirar un robot crea déficit. Trasladarlo a una tarea ocupada crea déficit y exceso, mientras que activar un robot sobre una tarea ocupada crea exceso. El ahorro de coste no supera $\kappa_{\max}<\lambda$, por lo que ninguna desviación mejora la utilidad. Los Nash son exactamente las $P(N,K)=N!/(N-K)!$ asignaciones factibles. Como $D=E=0$ en ese conjunto, los máximos globales del potencial coinciden con las asignaciones de coste mínimo.

Cada revisión rentable aumenta estrictamente \(\Phi_\lambda\). Los cambios aceptados visitan perfiles distintos de un conjunto con \((K+1)^N\) elementos, lo que acota la trayectoria por \((K+1)^N-1\) cambios. Bajo activación justa y ocupación coherente, el perfil terminal es Nash y satisface \(\boldsymbol z_0=\boldsymbol0\). El dominio de este resultado es el juego lógico con ocupación coherente; los modelos de red y planta se analizan por separado.

\subsection{Eficiencia y precios}
\label{app:proof-sp0-poa}

Todo óptimo social es factible y, por tanto, Nash. El precio de estabilidad es uno. No existe una cota uniforme para el peor Nash. En la instancia
\[
\boldsymbol\kappa_\epsilon=
\begin{bmatrix}0&\epsilon\\ \epsilon&1\end{bmatrix},
\qquad 0<\epsilon<\tfrac12,
\]
el perfil diagonal es Nash y cuesta $1$, mientras el cruzado cuesta $2\epsilon$. La razón $1/(2\epsilon)$ diverge cuando $\epsilon\downarrow0$.

\paragraph{Complementariedad $\varepsilon$.}
Defínase $v_i=\max_j(\beta_{ij}-p_j)$. Para cualquier asignación $\pi$ y una permutación $\sigma$ que satisface complementariedad $\varepsilon$ con precios compartidos,
\begin{align*}
 \sum_i\beta_{i,\pi(i)}&\leq\sum_iv_i+\sum_jp_j,\\
 \sum_i\beta_{i,\sigma(i)}&\geq\sum_iv_i+\sum_jp_j-N\varepsilon.
\end{align*}
Con $\beta=-\kappa$, se obtiene $C(\sigma)\leq C^\star+N\varepsilon$. Para costes enteros y $N\varepsilon<1$, la brecha debe ser cero.

\subsection{Precios locales, contracción espectral e imposibilidad}
\label{app:proof-sp0-local-prices}

Sea $\kappa^{\mathrm{pad}}$ la matriz cuadrada aumentada con las $N-K$ tareas ficticias de coste nulo introducidas en el Anexo~\ref{app:proof-sp0-complexity}. La complementariedad local
\[
 \kappa^{\mathrm{pad}}_{i,\sigma_t(i)}+\widehat\pi^t_{i,\sigma_t(i)}
 \leq\min_j\{\kappa^{\mathrm{pad}}_{ij}+\widehat\pi^t_{ij}\}+\varepsilon
\]
y $|\widehat\pi_{ij}-\bar\pi_j|\leq e_\pi^t$ implican complementariedad $(\varepsilon+2e_\pi^t)$ respecto al precio promedio. Al sumar sobre robots, los precios se cancelan porque $\sigma_t$ y una permutación óptima recorren una vez cada tarea. Se obtiene~\eqref{eq:sp0-local-price-gap}.

Defínase $\mathcal D_0:=\lVert[\widehat{\boldsymbol\pi}_i^0-\bar{\boldsymbol\pi}^0]_{i=1}^N\rVert_F$, el desacuerdo inicial apilado, que acota $\max_i\lVert\widehat{\boldsymbol\pi}_i^0-\bar{\boldsymbol\pi}^0\rVert_\infty$. Para $W=I-\alpha L$ simétrica, grafo conexo y $0<\alpha<2/\lambda_N(L)$, restringir $W$ al subespacio ortogonal a $\mathbf1$ da $e_\pi^t\leq q^t\mathcal D_0$, con $q=\max_{r\geq2}|1-\alpha\lambda_r(L)|<1$. Resolver $q^t\mathcal D_0\leq\eta$ produce $t_\eta=\lceil\log(\mathcal D_0/\eta)/(-\log q)\rceil$ y el coste suficiente $2|\mathcal E_C|Nt_\eta$. Esta cota cuenta transmisiones escalares y usa la contracción modal de peor caso; el empaquetado de valores o una contracción observada mayor reducirían el coste real.

Para la imposibilidad, considérense dos robots desconectados y dos tareas. El robot 1 observa la misma fila $(0,a)$, $0<a<1$, en dos instancias. Si la fila remota es $(0, 1)$, el óptimo le asigna la tarea 2. Si es $(1, 0)$, le asigna la tarea 1. Como su información local es idéntica, una decisión determinista coincide en ambos casos y falla en uno. Una red desconectada no puede garantizar optimalidad global para toda matriz de costes.
